\section{Sequential Neural Posterior Estimation}

Cuando hacemos SBI, en general tenemos que hacer muchas simulaciones para entrenar el modelo, la elección de prior es importante para que el proceso sea eficiente. Si el prior es muy amplio, la mayoría de las simulaciones van a caer en regiones de baja probabilidad, lo que hace que el entrenamiento sea ineficiente. Por eso, una idea natural es ir adaptando el prior a medida que vamos aprendiendo sobre el posterior, enfocándonos cada vez más en las regiones relevantes. La idea general de los métodos secuenciales de NPE (SNPE) se basan en esto: en vez de muestrear siempre desde el prior $p(\Theta)$, vamos adaptando de dónde sacamos los parámetros. En cada iteración usamos una \textit{proposal} $\tilde{p}(\Theta)$ que típicamente se parece al posterior que llevamos hasta el momento, de modo que concentramos las simulaciones en las regiones que realmente importan para explicar la observación $x_0$. Esto hace el proceso mucho más eficiente. Pero acá aparece un problema importante. Si entrenamos el modelo igual que antes, en realidad no estamos aprendiendo el posterior verdadero, sino uno sesgado por la proposal:
\[
q_\Phi^{\text{SNPE}}(\Theta \mid x) \approx \frac{1}{Z} p(x \mid \Theta)\tilde{p}(\Theta).
\]

donde $Z$ es una constante de normalización. Entonces, todos los métodos SNPE básicamente giran en torno a cómo corregir este sesgo introducido por usar una proposal distinta del prior. Ver \cite{greenberg2019} para una revisión de los distintos métodos SNPE (SNPE-A \cite{papamakarios2018fastepsilon}, SNPE-B y SNPE-C  o también llamado Automatic Posterior Transformation APT \cite{greenberg2019}).

Falta por investigar...

\subsection{Truncated Sequential Neural Posterior Estimation (TSNPE)}
Un enfoque alternativo a los métodos secuenciales tradicionales es el uso de un prior truncado \cite{deistler2022}. En lugar de simular parámetros en todo el espacio del prior original $p(\Theta)$, definimos una región de soporte $M$ basada en la posterior aproximada de la ronda actual, donde la densidad supera un umbral $\tau$. La propuesta truncada $\tilde{p}(\Theta)$ se define como:

\begin{equation}
\tilde{p}(\Theta) = \frac{1}{Z} p(\Theta) \cdot 1_{\Theta \in M}
\label{eq:truncated_prior}
\end{equation}

Donde $1_{\Theta \in M}$ es la función que vale $1$ si $\Theta \in M$ y $0$ en caso contrario. La constante de normalización $Z$ está dada por:

\begin{equation}
Z = \int_M p(\Theta) d\Theta
\label{eq:normalization_constant}
\end{equation}

En TSNPE, los datos se generan simulando a partir de la propuesta truncada, de modo que la distribución conjunta de entrenamiento pasa a ser $\tilde{p}(x,\Theta) = \tilde{p}(\Theta)p(x|\Theta)$. Siguiendo la ecuación \eqref{eq:loss_npe}, la función de pérdida de máxima verosimilitud bajo esta nueva propuesta es:

\begin{equation}
\mathcal{L}_{\text{TSNPE}} = - \mathbb{E}_{x,\Theta \sim \tilde{p}(x,\Theta)} \left[ \ln q_{\Phi}^{\text{NPE}}(\Theta|x) \right]
\label{eq:loss_tsnpe_expectation}
\end{equation}

Expandiendo la esperanza de forma explícita sobre el soporte restringido $M$, obtenemos:

\begin{equation}
\mathcal{L}_{\text{TSNPE}} = -\int_M \tilde{p}(\Theta) \left( \int p(x|\Theta) \ln q_{\Phi}^{\text{NPE}}(\Theta|x) dx \right) d\Theta
\label{eq:loss_integrals}
\end{equation}

Aplicando el Teorema de Bayes para sustituir la verosimilitud del simulador, sabemos que $p(x|\Theta) = \frac{p(x,\Theta)}{p(\Theta)}$. Al sustituir esta relación y la definición del prior truncado de la ecuación \eqref{eq:truncated_prior} dentro de la integral, se obtiene:

\begin{align}
\mathcal{L}_{\text{TSNPE}} &= -\int_M \left( \frac{1}{Z} p(\Theta) \right) \left( \int \frac{p(x,\Theta)}{p(\Theta)} \ln q_{\Phi}^{\text{NPE}}(\Theta|x) dx \right) d\Theta \\
&= -\frac{1}{Z} \int_M \int p(x,\Theta) \ln q_{\Phi}^{\text{NPE}}(\Theta|x) dx \, d\Theta \\
&= - \frac{1}{Z} \mathbb{E}_{x,\Theta \sim p(x,\Theta)} \left[ \ln q_{\Phi}^{\text{NPE}}(\Theta|x) \right] \cdot 1_{\Theta \in M}
\end{align}

Dado que la función de pérdida resultante es idéntica a la ecuación \eqref{eq:loss_npe} salvo por la constante multiplicativa $\frac{1}{Z}$ y la restricción del soporte a $M$, el gradiente con respecto a los parámetros de la red $\nabla_{\Phi} \mathcal{L}_{\text{TSNPE}}$ apunta exactamente en la misma dirección. Por lo tanto, no se requiere modificar analíticamente la función de pérdida, permitiendo utilizar el algoritmo de entrenamiento estándar de la primera ronda.